% main.tex  (FP2)
% Fixed Points II: Fixed Points in a 10D Modal Model
% Uses: series.sty (v3) and refs.bib
\documentclass[11pt]{article}
\usepackage{series}

% ---------------------------
% Paper metadata (override series defaults)
% ---------------------------
\renewcommand{\PartRoman}{II}
\hypersetup{
  unicode=true,
  pdftitle={Fixed Points II: Fixed Points in a 10D Modal Model},
  pdfauthor={Peter Nero}
}

% ---------------------------
% Local (paper) notation
% ---------------------------
\newcommand{\Mten}{\ensuremath{M_{10}}}
\newcommand{\Yfour}{\ensuremath{Y^{4}}}
\newcommand{\Xsix}{\ensuremath{X^{6}}}


\newcommand{\Aop}{A}

\newcommand{\HoneF}{\ensuremath{H^{1}_{F}}}
\newcommand{\Ltwo}{\ensuremath{L^{2}}}

\newcommand{\Lcoh}{L_{\mathrm{coh}}}

\newcommand{\kcen}{\kappa_{\mathrm{cen}}}
\newcommand{\Nil}{\mathrm{Nil}}
\newcommand{\Nilthree}{\mathrm{Nil}_{3}}

% For headings/bookmarks if needed
% series.sty already defines \FPpdf; keep a fallback just in case.
\providecommand{\FPpdf}[2]{\texorpdfstring{#1}{#2}}

\title{Fixed Points II: Projected Fixed Points and Equilibria in a 10D Modal Model}
\author{Peter Nero}
\date{July 2026\\Version 5}

\begin{document}
\maketitle
\seriestagline

\begin{abstract}
We apply the canonical FP--I machinery to a ten-dimensional control setting
$\Mten=\Yfour\times\Xsix$. The compact six-manifold $\Xsix$ carries three
compatible vertical structures represented by strongly commuting nonnegative
self-adjoint operators. Overlap is allowed; nesting requires supplied inclusion
maps and is not inferred from ranks $1<2<3$. In the q79 realization the
$1<2<3$ flag acts on a separate lane tensor factor, not inside an irreducible
HYM gauge bundle, while the shared circle is a separate flat line factor and
is not counted as a seventh internal product dimension. We prove existence of
projected time--$\tau$ fixed points by applying the canonical FP--I
Schauder/Darbo gates to this ten-dimensional operator, and prove coherent
uniqueness under base coercivity or
strong monotonicity. A projected fixed point is promoted to a full equilibrium
only under a strict Lyapunov identity. Fiber gaps control only the noncoherent
$Q$ sector and are never used as coherent damping.
\end{abstract}

\section*{Revision note for version 5}
\addcontentsline{toc}{section}{Revision note for version 5}
\begin{description}
\item[Supersedes.] \emph{Fixed Points II: Projected Fixed Points and
Equilibria in a 10D Modal Model}, version~4.
\item[Reason.] Version~4 gives the corrected ten-dimensional application, but
the relation between complement damping, coherent contraction, and
equilibrium promotion remained too compressed for a standalone reading.
\item[Resolution.] Version~5 adds a paper-specific reading guide and examples
that separate those three mechanisms, explains the lane factor and shared
line without adding dimensions, and keeps the generic fixed-point theorems
owned by Fixed Points~I. Publication description, revision history, and
computational provenance are now separate.
\item[Retained result.] The version~4 Schauder/Darbo application, conditional
coherent uniqueness, and strict-Lyapunov equilibrium gate are unchanged.
\item[Remaining boundary.] Physical q79 HYM endpoints, action, Hessian,
finite reduction, Lorentzian time, and observable interpretation remain open
downstream inputs.
\end{description}

\section*{Revision note for version 4}
\addcontentsline{toc}{section}{Revision note for version 4}
\begin{description}
\item[Supersedes.] \emph{Fixed Points II: Projected Fixed Points and
Equilibria in a 10D Modal Model}, version~3.
\item[Reason.] The version~3 correction fixed the ten-dimensional control
geometry but still described the q79 $1<2<3$ carrier too loosely. The exact
projective-module theorem excludes placing a nontrivial parallel flag inside
an irreducible stable HYM factor and distinguishes post-projection finite
algebra from a physical Galerkin subspace.
\item[Resolution.] Version~4 retains the corrected FP theorem and places the
rank flag on an external lane tensor factor, with the shared differential line
as a separate flat scalar factor. It also records that the 27-state algebra is
post-projection data and that the six-coordinate strain carrier is a nonlinear
quotient shadow rather than a linear subspace of the physical HYM carrier.
Generic existence and equilibrium-promotion results remain owned by FP--I;
this paper states their ten-dimensional application corollaries so that it is
standalone without duplicating the generic proofs.
\item[Retained result.] Schauder/Darbo existence and conditional coherent
uniqueness survive in the corrected ten-dimensional control realization.
\item[Remaining boundary.] The selected visible/hidden q79 HYM endpoints,
physical action and Hessian, finite invariant subspace or Feshbach execution,
physical time, and Lorentzian dynamics require independent source and
completion theorems.
\end{description}

\section*{How to read this application}
\addcontentsline{toc}{section}{How to read this application}
FP--I proves the generic analytic machinery.  This paper asks what that
machinery becomes in a ten-dimensional control model with one compact
six-dimensional internal space and three compatible vertical operators.  Its
formal statements are therefore application corollaries and
model-specialized estimates, not replacement copies of the FP--I theorems.

\paragraph{The central picture in plain language.}
There are two independent stabilization questions.  First, does the vertical
operator suppress every component outside the joint harmonic sector?  This is
controlled by the fiber gaps and concerns $\Qcoh\Psi$.  Second, does motion
inside the joint harmonic sector contract to one state?  The vertical
operators vanish there, so this requires a base Poincar\'e gap or strong
monotonicity of the coherent nonlinearity.  Confusing these mechanisms would
make a fiber gap appear to prove uniqueness when it only damps discarded
internal modes.

\paragraph{Argument map.}
Section~1 states the ten-dimensional control scope.  Section~2 places the
three operators on the same internal Hilbert space, constructs their joint
projector, and records the q79 lane/shared-line type boundary.  Section~3 then
separates the $\Pcoh$ and $\Qcoh$ evolutions, promotes a projected return only
under a strict Lyapunov identity, gives the two coherent contraction routes,
and transfers the compact-base and noncompact-base existence gates from
FP--I.  Section~4 supplies examples of the operator hypotheses.  The
appendices provide the specialized contraction and damping calculations.

\paragraph{What this paper does not select.}
The notation $\Mten=\Yfour\times\Xsix$ is a control realization, not a proof
that physical spacetime is this Riemannian product.  Nor does the paper derive
the q79 HYM bundle, identify the three operators with the $1<2<3$ lane
projectors, or turn the post-projection 27-state algebra into a continuum
Galerkin basis.  Those objects must enter through separate source and
intertwining theorems before the present estimates become physical
predictions.

\tableofcontents

% ============================================================
\section{Introduction and scope}
We consider $\Mten=\Yfour\times\Xsix$, where $(\Yfour,g_Y)$ is a complete
Riemannian analytic/control manifold and $\Xsix$ is compact. The Riemannian base
may model a Cauchy-slice control geometry or Euclideanized problem, but is not
identified here with physical Lorentzian spacetime. The internal space carries
three vertical structures indexed by $n=1,2,3$, each equipped with a
\emph{nonnegative} self-adjoint vertical operator and a uniform spectral gap.
We adopt throughout the corrected FP--I series conventions:
\begin{itemize}
\item Laplacians are \emph{nonnegative} ($\Delta\ge 0$), and dissipative flows are written
      $\partial_t\Psi=-\Aop\Psi-N(\Psi)$.
\item The coherent projector $\Pcoh$ is the joint fiber-harmonic projector, and
      $\Qcoh:=\Id-\Pcoh$.
\item Semigroup smoothing is recorded in the global safe form
      $\|\Aop^{1/2}e^{-t\Aop}\Qcoh\|\lesssim (1+t^{-1/2})e^{-\lambda t}$
      (cf.\ analytic semigroup bounds in \cite{Amann1995,Henry1981}).
\end{itemize}
This alignment is what allows FP--III and later papers to reuse damping/smoothing
constants without sign drift.

The parameter $t$ in the parabolic flow and the step $\tau$ below are
stabilization parameters. A physical Lorentzian evolution $U(t_2,t_1)$ and any
relation to this control flow require a separate theorem.

Our goals are:
(i) existence of coherent fixed points via Schauder (compact base) or Darbo (noncompact base with confinement);
(ii) uniqueness via a Fundamental Contractivity Condition (FCC) on the coherent sector;
(iii) explicit constants in standard geometries.
General dynamical-systems context may be compared with \cite{Temam1997}.

We also cite here FP--I as the foundational reference for the projection/fixed-point strategy:
see \cite{Nero2025FP1}.

% ============================================================
\section{Geometry and functional setup}

\subsection{Internal bundles and harmonic alignment}
Let $\Xsix$ be a compact six-manifold. On the common internal Hilbert space
$L^2(\Xsix)$ let $A_n(y)$, $n=1,2,3$, denote the nonnegative self-adjoint
operators associated with the selected vertical structures. These structures
may overlap. They are nested only if a concrete realization supplies bundle
injections or operator intertwiners; rank labels alone do not do so. No product
decomposition into three disjoint fibers is assumed. The shared central circle
is encoded by a common circle cycle, foliation, or $U(1)$ line-bundle
connection on $\Xsix$. Its line-bundle data do not add the dimension of the
bundle total space to $\dim\Xsix$.

\begin{remark}[Current q79 carrier and type boundary]
The current q79 preprojection architecture is
\[
\mathcal E_{\rm pre}
=\Gamma(E_{\rm HYM})\widehat\otimes H_{\rm lane}
\widehat\otimes L_{\rm shared}.
\]
The projectors
$p_1=\operatorname{diag}(1,0,0)$,
$p_2=\operatorname{diag}(1,1,0)$, and $p_3=I_3$ act only on
$H_{\rm lane}$. They encode the relative $1<2<3$ lanes without reducing the
holonomy of an irreducible stable HYM gauge factor. The common differential
line $L_{\rm shared}$ is a separate flat scalar factor and is not identified
with the curved HYM bundle. This placement is forced by the scalar commutant
of an irreducible stable HYM factor \cite{NeroQMSP2026}.

Two further type distinctions are essential. The accepted 27-state finite
algebra is post-projection source data, not a rank-27 Galerkin subspace of the
rank-102 physical deformation carrier. Likewise, the six real strain
coordinates are an orientation-forgetting nonlinear quotient shadow of
spectral data; there is no corresponding equivariant linear rank-six
subspace of $\operatorname{Herm}(3)$. The fixed-point theorems below use only
the declared operators and do not promote either finite object to the
physical continuum Hessian.
\end{remark}

We assume bounded internal geometry and smooth bounded dependence on
$y\in\Yfour$ as in FP--I. A concrete MTT realization must specify $\Xsix$, the
three operators, their domains, and the shared-circle action.

\begin{assumption}[Strong commutation of vertical operators]\label{ass:strong-comm}
For every $y$, the spectral measures of $A_1(y),A_2(y),A_3(y)$ commute. Their
common domain is dense, their weighted form sum is closed, and the commutation
and domain bounds are uniform in $y$. Equivalently, the three operators admit a
joint spectral calculus on the common internal Hilbert space.
\end{assumption}

\subsection{Fiber Laplacians and gaps}
Write $A_n(y)$ for the \emph{nonnegative} vertical Laplace-type operator
(or the vertical part of the spinor Laplacian in the spinor option). Each $A_n(y)$
has discrete spectrum with $0$ the harmonic eigenvalue. We assume a uniform gap
\begin{equation}\label{eq:gap}
\lamstar := \inf_{y\in \Yfour}\ \min_{1\le n\le 3}\ \lambda_1\!\big(A_n(y)\big) \;>\;0,
\end{equation}
and constant harmonic rank in $y$.

\paragraph{Explicit gaps in standard fibers.}
For $S^1_\ell$ (length $\ell$), $\lambda_1=(2\pi/\ell)^2$.
For a flat torus $T^m$ with side lengths $L_j$, $\lambda_1=(2\pi)^2\min_j L_j^{-2}$.
For the round $S^3$ (radius $1$), $\lambda_1=3$ (see, e.g., \cite{Chavel1984}).
Nilmanifold factors with bounded geometry admit a uniform positive lower bound on $\lambda_1$
once the metric class is fixed (bounded geometry prevents collapse).

\subsection{Fields, norms, and (optional) spinors}
We work with complex scalar fields by default; all results extend to sections of a fixed
Hermitian bundle with compatible connection.

\paragraph{Anisotropic norm.}
Let
\[
\HoneF := \bigcap_{n=1}^3\mathrm{dom}(A_n^{1/2}),
\qquad
\|\Psi\|_{\HoneF}^{2}:=\|\Psi\|_{\Ltwo}^{2}
+\sum_{n=1}^3\|A_n^{1/2}\Psi\|_{\Ltwo}^{2},
\]
with the natural common form domain. This is the joint vertical anisotropic
norm for the possibly overlapping structures. Under strong commutation it is
equivalent to the form norm of the weighted vertical sum when the weights
$\kappa_n$ are bounded above and below away from zero.

\subsection{Joint harmonic projector}
Let $\Pi_n(y)$ be the $\Ltwo$--orthogonal projector onto $\ker A_n(y)$ and define
the joint projector fiberwise by
\[
\Pcoh(y) := \Pi_1(y)\,\Pi_2(y)\,\Pi_3(y).
\]
We view $\Pcoh$ as the induced projector on $\Ltwo(\Mten)$ acting fiberwise.

\begin{lemma}[Commutation and range]\label{lem:comm-range}
Under Assumption~\ref{ass:strong-comm}, the spectral projectors commute. Hence
\[
\Pcoh(y)=\Pi_1(y)\Pi_2(y)\Pi_3(y),
\qquad
\Ran \Pcoh(y)=\bigcap_{n=1}^{3}\ker A_n(y).
\]
\end{lemma}

\paragraph{Geometric meaning of the joint projector.}
Each $\Pi_n$ asks whether a field is harmonic for one selected vertical
operator.  Their product retains only fields that answer yes to all three
questions at once.  Strong commutation is what makes this simultaneous test a
genuine orthogonal projector rather than an order-dependent sequence of
filters.  Overlap of the underlying vertical structures is harmless at this
operator level; literal nesting still requires independently supplied maps.

\begin{corollary}[10D joint-projector Sobolev regularity]\label{lem:Pcoh-H1}
Under the gap condition \eqref{eq:gap} and bounded fiber geometry, the joint projector
$\Pcoh:H^{1}\to H^{1}$ is bounded with
\[
\|\Pcoh\|_{H^{1}\to H^{1}}\le C_{\Pi},
\]
and $\Ran \Pcoh$ is closed in $H^{1}$.
\end{corollary}

This is the commuting three-projector specialization of the FP--I Sobolev
projector theorem \cite{Nero2025FP1}; the Riesz--Dunford and elliptic
regularity hypotheses are stated here explicitly.

\begin{remark}[Reference for Riesz projector bounds]
Uniform control of fiberwise Riesz projectors on Sobolev scales follows from the Riesz--Dunford formula
and parameter-dependent elliptic regularity; see \cite{Kato1976,Hebey2000}.
\end{remark}

% ============================================================
\section{Flow, projection, and fixed points}

\subsection{Parabolic generator and smoothing}
Let $\Delta_Y\ge0$ be the Riemannian Laplacian of the base control geometry
(or the Laplacian on a selected Riemannian Cauchy slice with declared boundary
conditions). It is not the Lorentzian d'Alembertian. We study
\begin{equation}\label{eq:flow}
\partial_t\Psi
= -\Bigl(\sum_{n=1}^{3}\kappa_n A_n + \varepsilon \Delta_{Y}\Bigr)\Psi
- N(\Psi),
\qquad
\kappa_n>0,\ \varepsilon\in[0,1].
\end{equation}
Denote the time--$\tau$ map by $\Phi_\tau$ and set
\[
\Qcoh := \Id-\Pcoh,
\qquad
\Aop := \sum_{n=1}^{3}\kappa_n A_n + \varepsilon \Delta_Y .
\]
Define rates
\[
\lambda_n := \inf_{y\in\Yfour}\lambda_1\!\big(A_n(y)\big)>0,
\qquad
\lambda_{\Aop} := \min_{1\le n\le 3}\kappa_n \lambda_n.
\]
Under strong commutation, $\lambda_{\Aop}$ is a lower bound for the vertical
form sum on $\Ran \Qcoh$.

\begin{assumption}[Projector commutation for the base-regularized model]\label{ass:projector-comm}
The common domain is preserved by $\Pcoh$ and $[\Aop,\Pcoh]=0$. If the vertical
operators depend on $y$, this includes control of the commutator
$[\Delta_Y,\Pcoh]$; it is not automatic from strong commutation of the
$A_n$. Models with nonzero commutator require explicit leakage terms and are
outside the exact invariant-sector theorem below.
\end{assumption}

By analytic semigroup theory (e.g.\ \cite{Amann1995,Henry1981}) and
Assumption~\ref{ass:projector-comm},
for $t>0$ one has
\begin{equation}\label{eq:semigroup}
\|\Aop^{1/2}e^{-t\Aop}\Qcoh\|_{\Ltwo\to\Ltwo}
\le C\,(1+t^{-1/2})e^{-\lambda_{\Aop}t},
\qquad
\|e^{-t\Aop}\Pcoh\|_{\Ltwo\to\Ltwo}\le 1.
\end{equation}
For two solutions let $w=\Psi_1-\Psi_2$. Duhamel's formula gives the
sector-correct estimate
\begin{align}
\|\Qcoh w(t)\|_{\HoneF}
&\le C(1+t^{-1/2})e^{-\lambda_{\Aop}t}
\|\Qcoh w(0)\|_{\Ltwo}\nonumber\\
&\quad+C\int_0^t(1+(t-s)^{-1/2})e^{-\lambda_{\Aop}(t-s)}
\|\Qcoh(N(\Psi_1(s))-N(\Psi_2(s)))\|_{\Ltwo}\,\dd s.
\label{eq:Q-duhamel}
\end{align}
No factor $e^{-\lambda_{\Aop}t}$ is asserted for $\Pcoh w$. On the invariant
coherent sector its dynamics are instead
\begin{equation}
\partial_t u=-\varepsilon\Delta_Yu-\Pcoh N(u),
\qquad u=\Pcoh u,
\label{eq:P-flow}
\end{equation}
and require base coercivity or coherent monotonicity for contraction.

\paragraph{Why the sector split matters.}
Equation~\eqref{eq:Q-duhamel} says that positive vertical modes are damped up
to nonlinear forcing.  Equation~\eqref{eq:P-flow} says that harmonic modes do
not feel that vertical gap at all.  Thus approaching the coherent sector and
selecting a unique point inside it are two different mathematical events.
The first can occur while an entire family of coherent states remains.

\begin{remark}[Base lower-order terms]
If the $\Qcoh$ equation includes first/zero-order terms with relative bound
$C_Y\ge0$, replace the noncoherent decay margin by
$\lambda_{\Aop}-C_Y$. This modification remains a $\Qcoh$-sector statement.
\end{remark}

\subsection{Coherence invariance and the projected map}
\begin{assumption}[Coherence invariance]\label{ass:coh-inv}
We assume Assumption~\ref{ass:projector-comm} and
$N(\Ran\Pcoh)\subset \Ran\Pcoh$.
\end{assumption}

We apply fixed-point principles to
\[
T_\tau := \Pcoh\circ \Phi_\tau:\ D\to \Ran\Pcoh\cap \Ltwo,
\]
where $D\subset \Ran\Pcoh\cap \Ltwo$ is a closed bounded set (specified in
\S\ref{subsec:exist}).

\begin{assumption}[Strict Lyapunov identity]\label{ass:lyapunov}
There is a functional $\mathcal E$ bounded below on the invariant set such that
every sufficiently regular trajectory satisfies
\[
\mathcal E(\Phi_tu)+c\int_0^t\|\partial_s\Phi_su\|_{\Ltwo}^2\,\dd s
=\mathcal E(u),\qquad c>0.
\]
For a gradient flow this follows from $N=\nabla V$ with the required domain and
regularity. It is not implied by semilinearity alone.
\end{assumption}

\begin{corollary}[10D projected-step equilibrium promotion]\label{prop:proj-eq}
Under Assumption~\ref{ass:coh-inv}, a point
$\Psi^\ast\in\Ran\Pcoh$ with $T_\tau(\Psi^\ast)=\Psi^\ast$ is a fixed point of
the time--$\tau$ stabilization map. If Assumption~\ref{ass:lyapunov} also holds,
then it is a genuine equilibrium:
$\Phi_t(\Psi^\ast)=\Psi^\ast$ for all $t\ge0$.
\end{corollary}

\begin{proof}
Apply the FP--I promotion theorem \cite{Nero2025FP1}. In the present
realization, coherence invariance gives
$T_\tau(\Psi^\ast)=\Phi_\tau(\Psi^\ast)=\Psi^\ast$. The Lyapunov identity on
$[0,\tau]$ then has equal endpoint energies, so its nonnegative dissipation
integral vanishes. Hence $\partial_t\Phi_t(\Psi^\ast)=0$ and the orbit is
stationary. Without the Lyapunov identity, a time--$\tau$ fixed point could be
a nonstationary periodic point.
\end{proof}

\paragraph{What equilibrium promotion establishes.}
The corollary closes a logical gap, not a dynamical one.  A projected
time--$\tau$ return becomes an equilibrium because strict energy loss forbids
a nonstationary loop with equal endpoint energy.  It does not prove that a
return exists or that it is unique; those are the separate existence and FCC
questions below.

\subsection{Uniqueness via the FCC (anisotropic norm)}
\begin{theorem}[Fundamental contractivity on the coherent sector]\label{thm:FCC}
Let
\[
\Lcoh := \mathrm{Lip}\bigl(\Pcoh N|_{\Ran \Pcoh}\bigr).
\]
Then for $\Psi_1,\Psi_2\in\Ran\Pcoh$,
\[
\|\Phi_\tau(\Psi_1)-\Phi_\tau(\Psi_2)\|_{\Ltwo}
\le e^{\Lcoh\tau}\|\Psi_1-\Psi_2\|_{\Ltwo},
\]
and since $\|\cdot\|_{\HoneF}=\|\cdot\|_{\Ltwo}$ on $\Ran\Pcoh$, the same holds in $\HoneF$.
Moreover, $T_\tau=\Pcoh\Phi_\tau$ is a Banach contraction on the declared
coherent phase space provided either:
\begin{enumerate}[label=\textup{(\alph*)}]
\item \textbf{Base diffusion:} $\varepsilon>0$ and the nonnegative
      $\Delta_Y$ has a Poincar\'e gap $\mu_Y>0$ on a specified invariant
      subspace (for example mean-zero functions or Dirichlet boundary data), with
      $\varepsilon\mu_Y>\Lcoh$; then
      \[
      \|T_\tau(\Psi_1)-T_\tau(\Psi_2)\|_{\Ltwo}\le e^{-(\varepsilon\mu_Y-\Lcoh)\tau}\|\Psi_1-\Psi_2\|_{\Ltwo}.
      \]
\item \textbf{Strong monotonicity:} there exists $\mu>0$ such that
      $\langle \Pcoh N(u)-\Pcoh N(v),u-v\rangle_{\Ltwo}\ge \mu\|u-v\|_{\Ltwo}^2$
      for all $u,v\in\Ran\Pcoh$; then
      \[
      \|T_\tau(\Psi_1)-T_\tau(\Psi_2)\|_{\Ltwo}\le e^{-\mu\tau}\|\Psi_1-\Psi_2\|_{\Ltwo}.
      \]
\end{enumerate}
\end{theorem}

\begin{remark}[Base zero mode]
On a compact boundaryless connected base, scalar constants lie in
$\ker\Delta_Y$. Therefore base diffusion gives no coherent contraction on the
full scalar space. One must remove/fix the mean, impose boundary conditions, or
use coherent strong monotonicity.
\end{remark}

\begin{remark}[Why fiber smoothing does not contract on $\Ran\Pcoh$]
On $\Ran\Pcoh$ one has $\nabla_F\Psi=0$, hence $\|\Psi\|_{\HoneF}=\|\Psi\|_{\Ltwo}$ and
$\Qcoh=0$. Thus contraction must come from \emph{base} dissipation (case (a)) or coherent monotonicity (case (b)),
not from the fiber semigroup estimate \eqref{eq:semigroup}.
\end{remark}

\paragraph{Interpretation of the FCC.}
The FCC is a uniqueness gate.  Its contraction factor compares two coherent
initial states after one stabilization step.  Base diffusion supplies that
factor only after its zero mode has been removed or fixed; strong
monotonicity supplies it directly from the nonlinear response.  Either route
can make Banach's theorem available, but neither is a consequence of the
internal spectral gap.

\subsection{Existence without contraction: Schauder and Darbo}\label{subsec:exist}

Uniqueness is useful when a contraction is available, but existence should not
be made to depend on it.  On a compact base, smoothing and Rellich compactness
lead to Schauder.  On a noncompact base, confinement replaces global
compactness by the compact-local/small-tail decomposition inherited from
FP--I.  The invariant set is part of either theorem and must be constructed in
the concrete model.

\begin{assumption}[Confinement for noncompact bases]\label{ass:conf}
If $\Yfour$ is noncompact, assume the energy includes a confining base potential $V(y)\to\infty$
as $r(y):=\dist_Y(y,y_0)\to\infty$, yielding uniform tail control on energy sublevels
(as in FP--I’s confinement hypothesis).
\end{assumption}

\begin{assumption}[Base smoothing for compactness]\label{ass:base-smooth}
For the Schauder/Darbo existence routes we assume \emph{base smoothing} is present, e.g.\
$\varepsilon>0$ in \eqref{eq:flow} (or one works with a base-regularized generator as in FP--I).
This ensures $\Phi_\tau$ maps bounded sets in $\Ltwo$ into sets bounded in $H^1$ on compact base regions,
so Rellich compactness applies.
\end{assumption}

\begin{lemma}[Coherent affine bound]\label{lem:affine}
Assume one of the two contraction hypotheses in Theorem~\ref{thm:FCC}, with
coherent contraction rate
\[
\rho_{\mathrm{coh}}=
\begin{cases}
\varepsilon\mu_Y-\Lcoh,&\text{base-gap case},\\
\mu,&\text{strong-monotonicity case}.
\end{cases}
\]
Set $q_{\mathrm{coh}}(\tau)=e^{-\rho_{\mathrm{coh}}\tau}$ and
$B_{\mathrm{coh}}(\tau)=\|T_\tau(0)\|_{\Ltwo}$. Then
\[
\|T_\tau(\Psi)\|_{\Ltwo}
\le q_{\mathrm{coh}}(\tau)\|\Psi\|_{\Ltwo}+B_{\mathrm{coh}}(\tau)
\]
for coherent $\Psi$. No $\Qcoh$-sector gap enters this estimate.
\end{lemma}

\begin{corollary}[Invariant/absorbing ball]\label{cor:ball}
If $q_{\mathrm{coh}}(\tau)<1$ in Lemma~\ref{lem:affine}, then
\[
D_R := \{\Psi\in \Ran\Pcoh\cap \Ltwo:\ \|\Psi\|_{\Ltwo}\le R\},
\qquad
R\ge \frac{B_{\mathrm{coh}}(\tau)}{1-q_{\mathrm{coh}}(\tau)},
\]
is forward invariant for $T_\tau$ and absorbing in the coherent $\Ltwo$ metric.
\end{corollary}

\begin{corollary}[Compact-base FP--I gate in the 10D model]\label{thm:schauder}
Assume $\Yfour$ is compact and Assumption~\ref{ass:base-smooth} holds. Then
$T_\tau=\Pcoh\Phi_\tau$ maps bounded $L^2$ sets into sets relatively compact
in $L^2$. If a nonempty $D\subset\Ran\Pcoh\cap\Ltwo$ is $L^2$-closed,
bounded, convex, and satisfies $T_\tau(D)\subset D$, then $T_\tau$ has a fixed
point in $D$.
\end{corollary}

\begin{lemma}[10D compact-local/small-tail reduction]\label{lem:condensing}
Assume Assumptions~\ref{ass:base-smooth} and \ref{ass:conf}. Let
$D_M$ be a forward-invariant energy sublevel that is bounded in $H^1$ and has
the uniform tail control of Assumption~\ref{ass:conf}. Let $\alpha(\cdot)$ be
the Kuratowski measure of noncompactness in $\Ltwo$. Then $\Phi_\tau$ is
\emph{Sadovski\u{\i}-condensing} on $D_M$: for every bounded
$E\subset D_M$ with $\alpha(E)>0$,
\[
\alpha(\Phi_\tau(E))<\alpha(E).
\]
\end{lemma}

\begin{proof}[Transfer of the FP--I mechanism]
Fix a bounded set $E$ in a forward-invariant energy sublevel. Let $\chi_R\in C_c^\infty(\Yfour)$ be a cutoff
equal to $1$ on $\{r\le R\}$ and supported on $\{r\le 2R\}$, and extend $\chi_R$ to $\Mten$ by $\chi_R\circ\pi$.
Define the output decomposition
\[
K_R(\Psi_0):=(\chi_R\circ\pi)\,\Phi_\tau(\Psi_0),
\qquad
S_R(\Psi_0):=(1-\chi_R\circ\pi)\,\Phi_\tau(\Psi_0).
\]
Confinement gives uniform small tails: $\sup_{\Psi_0\in E}\|S_R(\Psi_0)\|_{\Ltwo}\to 0$ as $R\to\infty$.
Base smoothing (Assumption~\ref{ass:base-smooth}) implies that for each fixed $R$, the set $K_R(E)$
is relatively compact in $\Ltwo$ (Rellich on the compact base region $\{r\le 2R\}$ plus smoothing in $H^1$).
Hence $\alpha(K_R(E))=0$, and by standard properties of $\alpha$,
\[
\alpha(\Phi_\tau(E))=\alpha(K_R(E)+S_R(E))\le \alpha(K_R(E)) + 2\sup_{\Psi_0\in E}\|S_R(\Psi_0)\|_{\Ltwo}
=2\sup_{\Psi_0\in E}\|S_R(\Psi_0)\|_{\Ltwo}.
\]
Choose $R$ so large that the right-hand side is $<\alpha(E)$, giving $\alpha(\Phi_\tau(E))<\alpha(E)$.
\end{proof}

\begin{corollary}[Noncompact-base FP--I gate in the 10D model]\label{thm:darbo}
Assume $\Yfour$ is noncompact and Assumptions~\ref{ass:conf} and \ref{ass:base-smooth} hold.
Let $D\subset \Ran\Pcoh\cap\Ltwo$ be nonempty, $L^2$-closed, bounded,
convex, and invariant under $T_\tau$. Since $\Pcoh$ is an $L^2$-orthogonal
projection of norm one, Lemma~\ref{lem:condensing} implies that $T_\tau$ is
condensing on $D$. Hence $T_\tau$ has a fixed point in $D$.
\end{corollary}

\begin{remark}[No full-map decay from the fiber gap]
The margin $\lambda_{\Aop}-C_Y-\Lip$ can sharpen estimates for the
$\Qcoh$ component. It cannot imply
$\alpha(\Phi_\tau(E))\le\theta(\tau)\alpha(E)$ with
$\theta(\tau)\to0$ for the full map unless the coherent component has its own
base or monotonicity contraction.
\end{remark}

\paragraph{What the existence gates contribute.}
The compact and condensing routes answer only whether $T_\tau$ has at least
one coherent fixed point in the declared invariant set.  Combined with the
strict Lyapunov identity, such a point is a full equilibrium.  Combined
further with the FCC, it is the unique equilibrium in the relevant coherent
phase space.  Keeping these three gates separate makes clear which conclusion
fails when compactness, Lyapunov strictness, or coherent contraction is absent.

\subsection{Damping estimates and effective rates}
Let $\Psi=\Pcoh\Psi+\Qcoh\Psi$ and write $\Psi^\perp:=\Qcoh\Psi$.
Assuming coherence invariance and that $N$ is $\Lip$--Lipschitz on the
invariant set considered, compare $\Psi$ with $\Pcoh\Psi$ to obtain
\[
\frac{d}{dt}\|\Psi^\perp\|_{\Ltwo}^{2}
\le -2\lambda_{\Aop}\|\Psi^\perp\|_{\Ltwo}^{2} + 2C_Y\|\Psi^\perp\|_{\Ltwo}^{2} + 2\Lip \|\Psi^\perp\|_{\Ltwo}^{2}.
\]
Hence the effective decay rate satisfies
\begin{equation}\label{eq:eta}
\eta \ge \lambda_{\Aop}-C_Y-\Lip.
\end{equation}
This governs noncoherent $\Ltwo$ decay only. Higher-norm control follows from
the sector-correct Duhamel estimate~\eqref{eq:Q-duhamel} under its regularity
hypotheses. It does not imply coherent or full $H^1$ decay.

% ============================================================
\section{Examples}

\subsection{Product of circles and tori}
Let $\Xsix=T^2_1\times T^2_2\times T^2_3$. Let $A_n$ contain the
nonnegative Laplacian in the $T_n^2$ directions, and let a common $U(1)$
connection encode the selected central-circle holonomy without adding another
product factor. For rectangular tori,
\[
\lambda_n
=\min\Bigl\{(2\pi/\ell_{u_n})^2,\ (2\pi/\ell_{v_n})^2\Bigr\}
>0,
\]
for the first nonconstant mode of that factor. These gaps control $\Qcoh$;
the FCC margin still requires base coercivity or coherent monotonicity.

\subsection{Sphere case (\FPpdf{$S^3$}{S3})}
If one selected vertical operator has round-$S^3$ leaves of radius one, its
first positive eigenvalue is $3$ (see \cite{Chavel1984}). This is an
operator-level gap example, not a claim that three disjoint $S^3$ factors fit
inside $\Xsix$ or that the coherent sector contracts.

\subsection{Nilmanifold auxiliary model}
On a compact six-manifold carrying a nilmanifold foliation or on an auxiliary
Lens--Nil model carrying a nilpotent vertical operator, a noncollapsing
bounded-geometry family has a
uniform positive first nonzero eigenvalue once its metric class is fixed. The
actual embedding, shared-circle action, and commutation with the other two
operators must be checked in the concrete model. This example does not identify
$L(3,1)\times\Nilthree$ with the selected q79/Fu--Yau compactification.

\subsection{Toward Physical Interpretation (programmatic note)}
While this paper is purely mathematical, the following schematic links guide later work:
\begin{itemize}
\item The three strongly commuting operators encode compatible, possibly
      overlapping modal structures on the same internal space.
\item The coherent sector $\Ran\Pcoh$ consists of fiberwise zero modes (only base dependence remains).
\item Spectral gaps $\lambda_n$ set decay rates of nonzero fiber modes; in KK heuristics this corresponds
      to mass scales $m_n\sim \lambda_n^{1/2}$ for excitations.
\item A projected step fixed point is an equilibrium only when the strict
      Lyapunov hypothesis of Proposition~\ref{prop:proj-eq} is verified.
\end{itemize}

% ============================================================
\section{Concluding remarks}
We specialized the FP--I framework to
$\Mten=\Yfour\times\Xsix$ with three strongly commuting vertical operators on
one compact internal space.  Their joint spectral projector gives a precise
coherent sector, vertical gaps damp its complement, and independent base or
monotonicity hypotheses control contraction inside it.  Compact-base
Schauder and noncompact-base Darbo gates provide alternative existence
mechanisms; a strict Lyapunov identity then distinguishes an actual
equilibrium from a merely periodic projected return.

The useful result is the resulting dependency chain:
\[
\begin{aligned}
\text{joint operator data}
&\Longrightarrow \text{coherent projector}
\Longrightarrow \text{existence},\\
\text{existence}
&\Longrightarrow \text{equilibrium promotion}
\Longrightarrow \text{conditional uniqueness}.
\end{aligned}
\]
No arrow may be skipped.  In particular, the fiber gap is not a coherent
contraction constant, the shared circle is not an extra internal dimension,
and the q79 lane flag is not a parallel subbundle of an irreducible HYM
factor.

What remains is a realization problem rather than another abstract
fixed-point lemma: a selected physical model must provide the q79 vertical
operators, their common domains and commutators, the physical action and
strict Lyapunov functional, and the relation between stabilization and
Lorentzian evolution.  Until those data are supplied, this paper establishes
a conditional ten-dimensional control theorem, not a derived physical
vacuum.

% ============================================================
\appendix

\section{Contraction on the coherent sector: details}
Throughout, let $T_\tau=\Pcoh\Phi_\tau$ act on the declared coherent phase
space and write $\Lcoh:=\mathrm{Lip}(\Pcoh N|_{\Ran\Pcoh})$.
Coherence invariance means $[\Aop,\Pcoh]=0$ and $\Pcoh N(\Ran\Pcoh)\subset\Ran\Pcoh$.

\begin{proof}[Proof of Theorem~\ref{thm:FCC} (baseline)]
If $\Psi_i(0)\in\Ran\Pcoh$, then $\Phi_t(\Psi_i)\in\Ran\Pcoh$ for all $t\ge 0$.
Set $u(t):=\Phi_t(\Psi_1)-\Phi_t(\Psi_2)\in\Ran\Pcoh$. Then
$\partial_t u = -\Pcoh\bigl(N(\Phi_t(\Psi_1))-N(\Phi_t(\Psi_2))\bigr)$, hence
$\frac{d}{dt}\|u\|_{\Ltwo}\le \Lcoh\|u\|_{\Ltwo}$ and Grönwall yields the bound.
\end{proof}

\begin{proof}[Proof of Theorem~\ref{thm:FCC}(a) (base diffusion)]
Assume $\varepsilon>0$ and the nonnegative $\Delta_Y$ has a Poincar\'e gap
$\mu_Y>0$ on the declared mean-zero or boundary-conditioned coherent
subspace. On that subspace the linear generator is $-\varepsilon\Delta_Y$,
so for $u(t)$ as above,
\[
\frac{d}{dt}\|u\|_{\Ltwo}^2 \le -2(\varepsilon\mu_Y-\Lcoh)\|u\|_{\Ltwo}^2,
\]
and Grönwall gives the stated contraction.
\end{proof}

\begin{proof}[Proof of Theorem~\ref{thm:FCC}(b) (strong monotonicity)]
Assume $\langle \Pcoh N(u)-\Pcoh N(v),u-v\rangle_{\Ltwo}\ge \mu\|u-v\|_{\Ltwo}^2$ on $\Ran\Pcoh$.
Then $\frac{d}{dt}\|u\|_{\Ltwo}^2\le -2\mu\|u\|_{\Ltwo}^2$ and Grönwall yields the contraction.
\end{proof}

\section{Spectral gap and damping inequality}
Let $\Psi=\Pcoh\Psi+\Qcoh\Psi$ and $\Psi^\perp:=\Qcoh\Psi$. For solutions of \eqref{eq:flow},
using that $\Aop\ge \lambda_{\Aop}I$ on $\Ran\Qcoh$ and that $N$ is $\Lip$--Lipschitz on the invariant set,
one obtains
\[
\frac{d}{dt}\|\Psi^\perp(t)\|_{\Ltwo}^2
\le -2(\lambda_{\Aop}-C_Y-\Lip)\|\Psi^\perp(t)\|_{\Ltwo}^2,
\]
which integrates to the exponential estimate consistent with \eqref{eq:eta}.

\section{Riesz projector bounds and base potentials}
For each $n$, let $A_n(y)\ge0$ be the nonnegative vertical operator defined in
Section~2 and fix a circle $\Gamma$ of radius $\lamstar/2$ centered at $0$
inside the spectral gap (uniform in $y$).
The fiberwise Riesz projector is
\[
\Pi_n(y) = \frac{1}{2\pi i}\int_\Gamma (z-A_n(y))^{-1}\,dz.
\]
Strong commutation gives the joint projector
\[
\Pcoh(y)=\Pi_1(y)\Pi_2(y)\Pi_3(y),
\qquad
\|\Pcoh\|_{\Ltwo\to\Ltwo}\le 1.
\]
Differentiating with respect to a base coordinate $y^\alpha$ gives
\[
\partial_\alpha \Pi_n(y)
=\frac{1}{2\pi i}\int_\Gamma (z-A_n(y))^{-1}\,(\partial_\alpha A_n(y))\,(z-A_n(y))^{-1}\,dz,
\]
hence
\[
\|\partial_\alpha \Pi_n(y)\|_{\Ltwo\to\Ltwo}
\le \frac{C}{(\lamstar)^2}\|\partial_\alpha A_n(y)\|_{H^2\to \Ltwo}.
\]
Under bounded internal geometry and $C^m$ base dependence,
$\|\partial_\alpha A_n(y)\|_{H^2\to \Ltwo}$ is uniformly bounded,
so $\|\partial_\alpha \Pcoh(y)\|_{\Ltwo\to\Ltwo}\le C'_\Pi$ and
\[
\|\Pcoh \Psi\|_{H^1}\le (1+C'_\Pi)\|\Psi\|_{H^1}.
\]
This is the standard mechanism behind Lemma~\ref{lem:Pcoh-H1}; see \cite{Kato1976,Hebey2000}.

\paragraph{Base potentials.}
Adding a nonnegative confining potential $V(y)$ improves tail coercivity on $\Yfour$ and underlies Darbo on noncompact bases,
but does not change the fiber operators $A_n(y)$ and hence does not change $\lamstar$ or $\lambda_{\Aop}$.

% ============================================================
\bibliographystyle{plain}
\bibliography{references.bib}

% BEGIN MTT MANAGED COMPUTATIONAL EVIDENCE
\section*{Computational Evidence and Reproducibility}
The numerical and machine-verifiable claims used by this paper are archived in the curated repository, \texttt{https://github.com/PeterNero/mtt-results-repro}. The mapped authority/result identifiers are \texttt{no result rows mapped}. Claim tiers in that capsule distinguish exact derivation, certified numerics, profile replay, conditional results, and open obligations.
% END MTT MANAGED COMPUTATIONAL EVIDENCE

\end{document}
